\documentclass[../main.tex]{subfiles}
\begin{document}
\section{Earnings-announcement drift}\label{sec:drift}

\subsection{Economic intuition}
Prices tend to adjust to earnings news over several weeks rather than all at once, a pattern known as post-earnings drift \citep{chan1996}. Measuring the surprise usually requires analyst forecasts, which our data do not include. \citet{brandt2008} show that the stock's own abnormal return around the announcement works as a surprise measure, and that it predicts returns over the following weeks.

\subsection{Signal and trading rules}
For each announcement we measure the stock's return over the reaction window, net of its market exposure:
\begin{equation}
  \text{EAR} \;=\; r_{\text{stock}} - \beta\, r_{\text{SPY}} .
\end{equation}
The reaction window runs from the last close before the announcement to the first close after it. For an after-close announcement, that is the announcement day's close to the next day's close. For a before-open announcement, it is the previous close to the announcement day's close. The beta $\beta$ is estimated by regressing the stock's daily returns on SPY's over the previous 60 trading days, capped between 0.5 and 2 (average 1.22).

The morning after the reaction window ends, we buy the stock at the opening auction if EAR is positive and short it if EAR is negative. Each position is 10\% of capital, hedged with an offsetting SPY position of $\beta$ times its value. We close both after 20 trading days at the official close. Positions in different stocks can overlap.

\subsection{Implementation}
The data cover 163 announcements across the 15 stocks: 104 in 2022--2023 and 59 in 2024. All trades are auction or closing-price trades. The engine is \texttt{engines/earnings\_events.cpp}.

\subsection{Results}
Table~\ref{tab:drift} reports the results. The average hedged return per event was 0.50\% in 2022--2023 and 1.55\% in 2024 (Figure~\ref{fig:events}). Neither average is statistically different from zero, with $t$-statistics of 0.59 and 1.32. At the account level the strategy earned 3.4\% a year over cash in 2022--2023 and 12.8\% in 2024.

\begin{table}[ht]\centering\small
\caption{Earnings-announcement drift, 20-day hold, beta-hedged, after costs.}\label{tab:drift}
\begin{tabular}{lrrrrrr}
\toprule
 & \multicolumn{3}{c}{Per event} & \multicolumn{3}{c}{Account (annualized)} \\
\cmidrule(lr){2-4}\cmidrule(lr){5-7}
Period & Events & Mean & $t$-stat & Return & Sharpe & Max drawdown \\
\midrule
2022--2023 (fitting) & 104 & 0.50\% & 0.59 & 3.4\% & 0.64 & $-6.2\%$ \\
2024 (out-of-sample) & 59 & 1.55\% & 1.32 & 12.8\% & 1.63 & $-6.1\%$ \\
\bottomrule
\end{tabular}
\end{table}

Larger reactions drift more. Restricting to announcements with $|\text{EAR}| > 2\%$ raises the average to 1.22\% per event in 2022--2023 (80 events) and 1.83\% in 2024 (42 events). The 2024 result is partly market exposure: although each position is hedged with a 60-day beta, the strategy's realized beta to SPY in 2024 was 0.31.

\Fig{figures/fig_events.png}{Average beta-hedged return per event, with 95\% intervals clustered by exit date, for the two earnings strategies.}{fig:events}

\subsection{Assessment}
The drift is in the direction the literature predicts and is larger in 2024. There are only about 50 to 60 events a year across 15 stocks, so the averages are noisy. The leftover market exposure in 2024 suggests the hedge needs a more stable beta estimate.

\ifSubfilesClassLoaded{% Bibliography.
\bigskip\noindent\emph{The reference details below should be checked against the published versions before citing.}\par\medskip
\begin{thebibliography}{99}
\setlength{\itemsep}{2pt}
\small
\bibitem[Bailey and L{\'o}pez de Prado(2014)]{bailey2014} Bailey, D. H., \& L{\'o}pez de Prado, M. (2014). The deflated Sharpe ratio. \emph{Journal of Portfolio Management}, 40(5).
\bibitem[Barber et~al.(2013)]{barber2013} Barber, B. M., De George, E. T., Lehavy, R., \& Trueman, B. (2013). The earnings announcement premium around the globe. \emph{Journal of Financial Economics}, 108(1).
\bibitem[Brandt et~al.(2008)]{brandt2008} Brandt, M. W., Kishore, R., Santa-Clara, P., \& Venkatachalam, M. (2008). Earnings announcements are full of surprises. Working paper.
\bibitem[Chan et~al.(1996)]{chan1996} Chan, L. K. C., Jegadeesh, N., \& Lakonishok, J. (1996). Momentum strategies. \emph{Journal of Finance}, 51(5).
\bibitem[Cremers and Weinbaum(2010)]{cremers2010} Cremers, M., \& Weinbaum, D. (2010). Deviations from put-call parity and stock return predictability. \emph{Journal of Financial and Quantitative Analysis}, 45(2).
\bibitem[DeMiguel et~al.(2009)]{demiguel2009} DeMiguel, V., Garlappi, L., \& Uppal, R. (2009). Optimal versus naive diversification: How inefficient is the 1/N portfolio strategy? \emph{Review of Financial Studies}, 22(5).
\bibitem[Easley et~al.(1998)]{easley1998} Easley, D., O'Hara, M., \& Srinivas, P. S. (1998). Option volume and stock prices: Evidence on where informed traders trade. \emph{Journal of Finance}, 53(2).
\bibitem[Fama and MacBeth(1973)]{famamacbeth1973} Fama, E. F., \& MacBeth, J. D. (1973). Risk, return, and equilibrium: Empirical tests. \emph{Journal of Political Economy}, 81(3).
\bibitem[Frazzini and Lamont(2007)]{frazzini2007} Frazzini, A., \& Lamont, O. A. (2007). The earnings announcement premium and trading volume. NBER Working Paper 13090.
\bibitem[Johnson and So(2012)]{johnsonso2012} Johnson, T. L., \& So, E. C. (2012). The option to stock volume ratio and future returns. \emph{Journal of Financial Economics}, 106(2).
\bibitem[Jorion(1986)]{jorion1986} Jorion, P. (1986). Bayes-Stein estimation for portfolio analysis. \emph{Journal of Financial and Quantitative Analysis}, 21(3).
\bibitem[Ledoit and Wolf(2004)]{ledoit2004} Ledoit, O., \& Wolf, M. (2004). A well-conditioned estimator for large-dimensional covariance matrices. \emph{Journal of Multivariate Analysis}, 88(2).
\bibitem[Markowitz(1952)]{markowitz1952} Markowitz, H. (1952). Portfolio selection. \emph{Journal of Finance}, 7(1).
\bibitem[Newey and West(1987)]{neweywest1987} Newey, W. K., \& West, K. D. (1987). A simple, positive semi-definite, heteroskedasticity and autocorrelation consistent covariance matrix. \emph{Econometrica}, 55(3).
\bibitem[Pan and Poteshman(2006)]{pan2006} Pan, J., \& Poteshman, A. M. (2006). The information in option volume for future stock prices. \emph{Review of Financial Studies}, 19(3).
\bibitem[Xing et~al.(2010)]{xing2010} Xing, Y., Zhang, X., \& Zhao, R. (2010). What does the individual option volatility smirk tell us about future equity returns? \emph{Journal of Financial and Quantitative Analysis}, 45(3).
\end{thebibliography}
}{}
\end{document}
